\documentclass{article}
\usepackage[utf8]{inputenc}
\usepackage{amsmath}
\usepackage{geometry}
\geometry{margin=2cm}
\begin{document}

\section*{Análise da Questão 155 — ENEM 2024 — Matemática}

\subsection*{Enunciado}
Um proprietário deseja instalar um sensor de presença para proteger seu imóvel. A área de cobertura do sensor é um setor circular definido pelo ponto \(S\), com raio \(R\) e ângulo de abertura \(\alpha\), ambos variando conforme o modelo do sensor. Cinco tipos de sensores estão disponíveis, com especificações:
\begin{itemize}
  \item Tipo I: \(\alpha = 15^\circ\), \(R = 20\,m\);
  \item Tipo II: \(\alpha = 30^\circ\), \(R = 22\,m\);
  \item Tipo III: \(\alpha = 40^\circ\), \(R = 12\,m\);
  \item Tipo IV: \(\alpha = 60^\circ\), \(R = 16\,m\);
  \item Tipo V: \(\alpha = 90^\circ\), \(R = 10\,m\).
\end{itemize}
O proprietário requer um sensor com área mínima de cobertura de 70 \(\mathrm{m}^2\), desejando o modelo mais barato possível (menor área que atenda a essa condição). Considere \(\pi \approx 3\).

\subsection*{Solução}
Aplicamos a fórmula da área do setor circular:
\[
A = \frac{\alpha}{360} \times \pi \times R^2.
\]
Calculando para cada sensor:

\begin{tabular}{c|c|c|c}
Tipo & $\alpha$ (°) & $R$ (m) & $A$ ($\mathrm{m}^2$) \\
\hline
I & 15 & 20 & $\frac{15}{360} \times 3 \times 400 = 50$ \\
II & 30 & 22 & $\frac{30}{360} \times 3 \times 484 = 121$ \\
III & 40 & 12 & $\frac{40}{360} \times 3 \times 144 = 48$ \\
IV & 60 & 16 & $\frac{60}{360} \times 3 \times 256 = 128$ \\
V & 90 & 10 & $\frac{90}{360} \times 3 \times 100 = 75$ \\
\end{tabular}

Sensores que atendem a pelo menos 70 m²: tipos II, IV e V.

\subsection*{Escolha final}
Entre os sensores que cobrem a área mínima, o tipo V apresenta a menor área de 75 m², o que também indica menor preço. Portanto, o proprietário deve adquirir o sensor do tipo \textbf{V}.

\subsection*{Discussão de alternativas incorretas}
\begin{itemize}
\item Tipo I e III possuem áreas inferiores ao mínimo, logo não atendem.
\item Tipo II e IV atendem, porém têm áreas maiores que o tipo V, o que pressupõe maior custo.
\end{itemize}

\subsection*{Perfil da questão}
A questão avalia a habilidade de cálculo da área de setores circulares, interpretação de figuras geométricas e tomada de decisão baseada em comparação quantitativa.

\subsection*{Considerações finais}
A resposta correta é a opção que representa o sensor do tipo \textbf{V}.

\end{document}